% Options for packages loaded elsewhere
\PassOptionsToPackage{unicode}{hyperref}
\PassOptionsToPackage{hyphens}{url}
\documentclass[
  11pt,a4paper]{article}
\usepackage{xcolor}
\usepackage[margin=1in]{geometry}
\usepackage{amsmath,amssymb}
\setcounter{secnumdepth}{-\maxdimen} % remove section numbering
\usepackage{iftex}
\ifPDFTeX
  \usepackage[T1]{fontenc}
  \usepackage[utf8]{inputenc}
  \usepackage{textcomp} % provide euro and other symbols
\else % if luatex or xetex
  \usepackage{unicode-math} % this also loads fontspec
  \defaultfontfeatures{Scale=MatchLowercase}
  \defaultfontfeatures[\rmfamily]{Ligatures=TeX,Scale=1}
\fi
\usepackage{lmodern}
\ifPDFTeX\else
  % xetex/luatex font selection
  \setmainfont[]{Latin Modern Roman}
\fi
% Use upquote if available, for straight quotes in verbatim environments
\IfFileExists{upquote.sty}{\usepackage{upquote}}{}
\IfFileExists{microtype.sty}{% use microtype if available
  \usepackage[]{microtype}
  \UseMicrotypeSet[protrusion]{basicmath} % disable protrusion for tt fonts
}{}
\makeatletter
\@ifundefined{KOMAClassName}{% if non-KOMA class
  \IfFileExists{parskip.sty}{%
    \usepackage{parskip}
  }{% else
    \setlength{\parindent}{0pt}
    \setlength{\parskip}{6pt plus 2pt minus 1pt}}
}{% if KOMA class
  \KOMAoptions{parskip=half}}
\makeatother
\usepackage{longtable,booktabs,array}
\newcounter{none} % for unnumbered tables
\usepackage{calc} % for calculating minipage widths
% Correct order of tables after \paragraph or \subparagraph
\usepackage{etoolbox}
\makeatletter
\patchcmd\longtable{\par}{\if@noskipsec\mbox{}\fi\par}{}{}
\makeatother
% Allow footnotes in longtable head/foot
\IfFileExists{footnotehyper.sty}{\usepackage{footnotehyper}}{\usepackage{footnote}}
\makesavenoteenv{longtable}
\setlength{\emergencystretch}{3em} % prevent overfull lines
\providecommand{\tightlist}{%
  \setlength{\itemsep}{0pt}\setlength{\parskip}{0pt}}
\usepackage{mathtools}
\usepackage{bm}
\usepackage{enumitem}
\usepackage{microtype}
\usepackage{booktabs}
\usepackage{array}
\usepackage{bookmark}
\IfFileExists{xurl.sty}{\usepackage{xurl}}{} % add URL line breaks if available
\urlstyle{same}
\hypersetup{
  pdftitle={ED Mobility Saturation Predicts Non-Newtonian Rheology A Fluid-Mechanical Extension of the Universal Mobility Law},
  pdfauthor={Allen Proxmire},
  hidelinks,
  pdfcreator={LaTeX via pandoc}}

\title{ED Mobility Saturation Predicts Non-Newtonian Rheology\\
\large A Fluid-Mechanical Extension of the Universal Mobility Law}
\author{Allen Proxmire}
\date{April 2026}

\begin{document}
\maketitle

\subsection{Abstract}\label{abstract}

We extend the empirically-validated Universal Mobility Law
\(D(c) = D_0(1 - c/c_\mathrm{max})^\beta\) --- established for ten
chemically distinct soft-matter systems with \(R^2 > 0.986\) --- from
concentration-driven diffusion in soft matter to fluid-mechanical
rheology. Treating the Universal Mobility Law as the empirical content
of Event Density's architectural principle P4 (mobility capacity bound),
we derive five form-FORCED predictions: (i) the Krieger-Dougherty
viscosity divergence
\(\mu(\rho) = \mu_0(1 - \rho/\rho_\mathrm{max})^{-\beta}\) near
suspension jamming; (ii) discontinuous shear-thickening /
shear-thickening-fluid (STF) divergence under strain-rate-driven
mobility saturation; (iii) Cross-class shear-thinning under monotone
shear-rate-dependent mobility; (iv) mixed regime thinning-to-thickening
transitions admitted by non-monotone constitutive choices for the
mobility function; and (v) Maxwell-class viscoelasticity arising as the
leading-order coarse-graining of ED's V5 cross-chain memory kernel under
exponential-decay ansatz. Of thirteen catalogued standard rheology
classes, five are form-FORCED at the architectural canon level, four are
compatible-INHERITED (form-derivable with material-specific shape
parameters not predicted from primitives), and four are non-canonical
(Oldroyd-B, Giesekus, FENE-P, Bingham yield-stress) requiring
constitutive structure beyond P4 + V5. Empirical comparison against the
Universal Mobility Law's ten-system distribution gives mean
\(\beta = 1.72 \pm 0.37\), range \([1.30, 2.49]\), all \(R^2 > 0.986\);
this is consistent within \(1\sigma\) at individual-system scatter level
with the canonical Event Density prediction \(\beta = 2.0\).
Mechanism-clustered \(\beta\) patterns --- cooperative networks (H-bond
viscosity, hydrodynamic crowding) cluster at \(\bar\beta \approx 2.31\);
steric / canonical (hard-sphere, polymer suspensions) at
\(\bar\beta \approx 1.76\); gradual / less-cooperative (polymer
entanglement, electrostatic) at \(\bar\beta \approx 1.38\) --- are
structurally coherent with the framework's cooperativity interpretation.
The fluid-mechanical extension unifies soft-matter mobility,
suspension-rheology divergence, shear-thickening fluids, polymer-melt
thinning, and Maxwell viscoelasticity under a single architectural
principle (P4 + V5) with canonical \(\beta = 2.0\) as the quantitative
anchor. Parameter values (\(\beta\), \(\rho_\mathrm{max}\),
\(\dot\gamma_\mathrm{max}\), Cross / Carreau shape parameters, V5
relaxation time \(\tau_R\)) are INHERITED per material-specific physics.
The result extends ED's empirical reach from the soft-matter regime
(Universal Mobility Law) into mainstream non-Newtonian fluid mechanics,
with no new free parameters introduced beyond those already present in
standard rheology models.

\begin{center}\rule{0.5\linewidth}{0.5pt}\end{center}

\subsection{1. Introduction}\label{introduction}

\subsubsection{1.1 The empirical universality problem in non-Newtonian
rheology}\label{the-empirical-universality-problem-in-non-newtonian-rheology}

Non-Newtonian rheology --- the study of fluids whose viscous response
deviates from the linear Newtonian \(\sigma = \mu \dot\gamma\) relation
--- encompasses a remarkably diverse set of empirical phenomena: dense
colloidal suspensions exhibit Krieger-Dougherty divergence near
close-packing; cornstarch and silica suspensions exhibit discontinuous
shear-thickening at critical shear rates; polymer melts and biological
fluids exhibit Cross / Carreau-class shear-thinning; viscoelastic gels
and polymer solutions exhibit Maxwell-class memory-kernel response. Each
phenomenon is, individually, well-characterized empirically. None has a
single unifying first-principles derivation in standard fluid mechanics;
instead, each is described by a domain-specific phenomenological model
(Krieger-Dougherty for suspensions, Cross or Carreau for polymer fluids,
Maxwell or Oldroyd-B for viscoelastic media) with separately-fit
material parameters.

The question this paper addresses is whether these distinct rheology
classes share a deeper architectural origin. We argue that they do, and
that the architectural principle is the \emph{mobility capacity bound}
P4 of Event Density's canonical PDE --- the same principle that, applied
to concentration in soft matter, generates the empirically-validated
Universal Mobility Law.

\subsubsection{1.2 The Universal Mobility Law as empirical
anchor}\label{the-universal-mobility-law-as-empirical-anchor}

The Universal Mobility Law (UMD) {[}Proxmire 2026{]} establishes that
concentration-dependent self-diffusion in crowded soft matter follows
the universal form

\[D(c) = D_0 \left(1 - \frac{c}{c_\mathrm{max}}\right)^\beta\]

across ten chemically distinct materials --- colloids (hard-sphere,
PMMA, Ludox silica, casein), proteins (BSA, lysozyme), polymers
(PEG-water, dextran), small molecules (sucrose-water, glycerol-water)
--- with \(R^2 > 0.986\) in every fit. The fitted exponent ranges
\(\beta \in [1.30, 2.49]\) with mean \(\bar\beta = 1.72 \pm 0.37\),
consistent within \(1\sigma\) with the canonical ED-architectural value
\(\beta = 2.0\).

The UMD is the empirical content of Event Density's principle P4
(mobility capacity bound) instantiated on a scalar density field. It is
not a phenomenological fit specific to any particular material; it is a
structural prediction that the same functional form must hold across
mechanistically distinct systems whose only shared feature is finite
packing capacity.

\subsubsection{1.3 The fluid-mechanical extension
question}\label{the-fluid-mechanical-extension-question}

This paper extends the UMD's empirical anchor from soft-matter
concentration-driven diffusion to fluid-mechanical rheology. The
structural template is identical: P4 specifies that the mobility
function \(M(\cdot)\) vanishes at a packing-class upper bound
\(\cdot_\mathrm{max}\) on its argument. By varying which physical
variable plays the role of the saturating argument --- fluid density
\(\rho\), strain-rate magnitude \(\dot\gamma\), or generalized state
variables --- and by combining P4 with V5 (cross-chain memory kernel for
time-dependent response), we recover a family of predicted rheology
classes.

The architectural commitment is fixed; the constitutive choice (which
argument saturates, what specific functional form \(M\) takes within the
form-FORCED class) is INHERITED per material physics. This pattern ---
\emph{form-FORCED, value-INHERITED} --- is the canonical methodology of
the Event Density program and is preserved here.

\subsubsection{1.4 Scope and
contributions}\label{scope-and-contributions}

The paper's substantive contributions:

\begin{enumerate}
\def\labelenumi{\arabic{enumi}.}
\item
  \textbf{Density-driven jamming → Krieger-Dougherty divergence}
  (Section 3): the universal viscosity-divergence form near suspension
  close-packing follows from P4 mobility saturation combined with the
  standard Stokes-Einstein-class mobility-viscosity inversion.
\item
  \textbf{Strain-rate extension → DST, STF, and Cross / Carreau
  shear-thinning} (Section 4): applying the same architectural template
  with \(\dot\gamma\) as the saturating argument yields shear-thickening
  / discontinuous shear-thickening (mobility vanishing at upper
  shear-rate bound) and Cross-class shear-thinning (monotone-decreasing
  mobility).
\item
  \textbf{Rheology-class catalogue} (Section 5): we classify thirteen
  standard rheology classes by their architectural status --- five
  form-FORCED, four compatible-INHERITED, four non-canonical.
\item
  \textbf{Empirical comparison} (Section 6): the UMD's ten-system
  \(\beta\) distribution provides quantitative empirical anchoring;
  canonical ED prediction \(\beta = 2.0\) is consistent within
  \(1\sigma\) of the empirical mean \(\bar\beta = 1.72 \pm 0.37\).
\item
  \textbf{V5 → Maxwell viscoelasticity} (Section 7): ED's cross-chain
  memory kernel, under exponential-decay ansatz, reduces to the Maxwell
  viscoelastic constitutive equation
  \(\tau_R \dot\sigma + \sigma = 2\mu S\), with \(\tau_R\) and amplitude
  INHERITED.
\item
  \textbf{Architectural synthesis} (Section 8): we show that the
  fluid-mechanical rheology landscape unifies under a single
  architectural framework (P4 + V5) with canonical \(\beta = 2.0\) as
  the quantitative anchor.
\end{enumerate}

The non-canonical classes --- Oldroyd-B, Giesekus, FENE-P, Bingham
yield-stress --- are flagged honestly as requiring constitutive
structure beyond the canonical P4 + V5 framework. They are real
empirical phenomena that ED's current canon does not natively absorb; we
discuss their architectural status and candidate canon-extension routes
in Section 9.

\subsubsection{1.5 Relationship to standard non-Newtonian
rheology}\label{relationship-to-standard-non-newtonian-rheology}

ED's prediction does not contradict standard non-Newtonian rheology; it
provides a structural origin for it. The Krieger-Dougherty form, the
Cross and Carreau models, and the Maxwell viscoelastic equation remain
valid as the canonical phenomenological descriptions. ED's contribution
is to derive their \emph{functional form} from a single architectural
principle, with material-specific parameters (the \(\beta\) exponent,
packing fractions, shape parameters, relaxation times) inherited per
system rather than postulated. The empirical match between the
ten-system UMD distribution and the canonical \(\beta = 2.0\) within
\(1\sigma\) is the substantive quantitative content of the paper.

\begin{center}\rule{0.5\linewidth}{0.5pt}\end{center}

\subsection{2. Architectural Background}\label{architectural-background}

\subsubsection{2.1 The Event Density canonical
PDE}\label{the-event-density-canonical-pde}

Event Density's canonical scalar partial differential equation,
established in the program's foundational papers, takes the form

\[\partial_t \rho \;=\; D \cdot F[\rho] \;+\; H \cdot v, \qquad D + H = 1, \qquad D, H \in [0, 1],\]

\[\dot v = \frac{1}{\tau}\bigl(\bar F(\rho) - \zeta v\bigr),\]

with the operator

\[F[\rho] = M(\rho) \nabla^2 \rho + M'(\rho) |\nabla\rho|^2 - P(\rho).\]

Three constitutive channels --- mobility (\(M\)), penalty (\(P\)), and
participation (\(v\)) --- are jointly necessary and sufficient under the
seven structural constraints C1--C7 of {[}Foundations 2025{]}. The
architectural canon abstracts these constraints to the seven principles
P1--P7 {[}Architectural Canon 2026{]}; the Universal Mobility Law is the
empirical signature of P4 (mobility capacity bound) operating on a
scalar field.

\subsubsection{2.2 P4: the mobility capacity
bound}\label{p4-the-mobility-capacity-bound}

Principle P4 specifies:

\[M(\rho_\mathrm{max}) = 0, \qquad M(\rho) > 0 \text{ for } \rho < \rho_\mathrm{max}.\]

The mobility function vanishes at a packing-class upper bound, and is
positive below. The simplest functional class admissible under the four
constitutive constraints (non-negativity, vanishing at capacity,
smoothness, dissipative structure) is the monomial form

\[M(\rho) = M_0 (1 - \rho/\rho_\mathrm{max})^\beta, \qquad \beta > 0.\]

The exponent \(\beta\) is INHERITED per material structure. The
canonical ED-architectural value is \(\beta = 2.0\), which arises as the
natural exponent under multiplicative-cooperativity counting (each unit
of jamming-cooperativity contributes one factor of
\((1 - \rho/\rho_\mathrm{max})\) to the mobility suppression).

\subsubsection{2.3 V5: cross-chain memory
kernel}\label{v5-cross-chain-memory-kernel}

Principle V5 (vacuum-induced cross-chain correlations), established in
arc-N N.2 of the Event Density program, specifies that same-sector
chains acquire correlated bandwidth content via vacuum-mediated
coupling. Translated to fluid-mechanical language, V5 introduces a
finite-memory kernel \(K_\mathrm{V5}(t)\) relating present stress to
past strain rate:

\[\sigma_\mathrm{V5}(t) = 2\mu \int_{-\infty}^t K_\mathrm{V5}(t - t') \, S(t') \, dt'.\]

The memory time \(\tau_R\) --- the characteristic decay scale of
\(K_\mathrm{V5}\) --- is INHERITED; it depends on molecular-relaxation
physics specific to the material. For polymer melts and solutions
\(\tau_R \in [10^{-3}, 10^0]\,\mathrm{s}\); for biological gels
\(\tau_R \in [10^{-2}, 10^1]\,\mathrm{s}\); for water
\(\tau_R \sim 10^{-12}\,\mathrm{s}\).

\subsubsection{2.4 The dual mobility-viscosity
mapping}\label{the-dual-mobility-viscosity-mapping}

Two distinct mobility-to-viscosity identifications are operative in
fluid mechanics, depending on the physical scenario:

\begin{itemize}
\item
  \textbf{Stokes-Einstein-class inversion} (concentration-mobility
  scenario): \(\mu \propto 1/M\). Used when \(M\) represents particle
  self-diffusion in a suspension and viscosity is the bulk viscoelastic
  response. Operative in dense suspensions near jamming (Section 3).
\item
  \textbf{NS-direct mapping} (velocity-mobility scenario):
  \(\mu = D \cdot M\). Used when \(M\) represents the
  kinematic-viscosity-class diffusivity for velocity gradients.
  Operative in polymer-melt-class shear-thinning (Section 4).
\end{itemize}

Both are admissible at the architectural level; the choice is determined
by which physical mobility the Event Density mobility function \(M\)
represents in the particular system. The Stokes-Einstein-class inversion
is imported from kinetic theory (not derived from ED); ED supplies the
saturation form of \(M\), kinetic theory supplies the inversion.

\subsubsection{2.5 The form-FORCED / value-INHERITED
methodology}\label{the-form-forced-value-inherited-methodology}

Throughout this paper we maintain the canonical Event Density
distinction:

\begin{itemize}
\tightlist
\item
  \textbf{Form-FORCED}: structural functional form derived from
  architectural principles. Universal across systems.
\item
  \textbf{Value-INHERITED}: specific numerical parameters within the
  form-FORCED class. Determined by material-specific physics, not
  predictable from primitives alone.
\end{itemize}

The Universal Mobility Law is form-FORCED via P4; the specific \(\beta\)
value for any given system is INHERITED. The canonical \(\beta = 2.0\)
is a \emph{predicted central value} under the architectural
multiplicative-cooperativity reading; empirical
\(\bar\beta = 1.72 \pm 0.37\) is consistent within \(1\sigma\).

\begin{center}\rule{0.5\linewidth}{0.5pt}\end{center}

\subsection{3. D1 --- Density-Driven
Jamming}\label{d1-density-driven-jamming}

\subsubsection{3.1 Setup: P4 mobility saturation in dense
suspensions}\label{setup-p4-mobility-saturation-in-dense-suspensions}

For a dense suspension of particles in a Newtonian carrier fluid, the
relevant mobility is the particle self-diffusion coefficient
\(M(\rho)\), where \(\rho\) is the suspension's volume fraction (or
equivalent packing density). The Universal Mobility Law fixes the
functional form:

\[M(\rho) = M_0 \left(1 - \frac{\rho}{\rho_\mathrm{max}}\right)^\beta, \qquad \beta > 0,\]

with \(M(\rho_\mathrm{max}) = 0\) corresponding to physical jamming at
close-packing.

\subsubsection{3.2 Stokes-Einstein-class mobility-viscosity
inversion}\label{stokes-einstein-class-mobility-viscosity-inversion}

For Brownian particles in moderate-to-dense suspensions, the
Stokes-Einstein relation in the dilute limit gives

\[D_\mathrm{tracer} = \frac{k_B T}{6\pi \mu a},\]

with \(a\) the particle radius. Inverting to express viscosity as a
function of mobility:

\[\mu \propto \frac{1}{M(\rho)}.\]

For finite-concentration suspensions the proportionality constant is set
by the dilute limit (\(\mu_\mathrm{susp}(0) = \mu_0\), the carrier-fluid
viscosity; \(M(0) = M_0\)):

\[\mu_\mathrm{susp}(\rho) = \mu_0 \cdot \frac{M_0}{M(\rho)}.\]

This is the standard Stokes-Einstein-class extension to moderately-dense
suspensions (Russel, Saville \& Schowalter 1989).

\subsubsection{3.3 Result: Krieger-Dougherty
form}\label{result-krieger-dougherty-form}

Substituting the UMD mobility form:

\[\boxed{\;\mu_\mathrm{susp}(\rho) = \mu_0 \left(1 - \frac{\rho}{\rho_\mathrm{max}}\right)^{-\beta}.\;}\]

This is the \textbf{Krieger-Dougherty form} (Krieger \& Dougherty 1959),
structurally identical to the standard empirical model for suspension
viscosity divergence near close-packing. The standard Krieger-Dougherty
model writes

\[\frac{\mu_\mathrm{susp}}{\mu_0} = \left(1 - \frac{\phi}{\phi_\mathrm{max}}\right)^{-[\eta]\phi_\mathrm{max}},\]

with intrinsic viscosity \([\eta] = 5/2\) for hard spheres and maximum
packing fraction \(\phi_\mathrm{max} \approx 0.64\) (random close
packing), giving exponent \([\eta]\phi_\mathrm{max} \approx 1.82\). Our
derived form maps via \(\rho \leftrightarrow \phi\),
\(\rho_\mathrm{max} \leftrightarrow \phi_\mathrm{max}\),
\(\beta \leftrightarrow [\eta]\phi_\mathrm{max}\).

\subsubsection{3.4 What is form-FORCED, what is INHERITED, what is
imported}\label{what-is-form-forced-what-is-inherited-what-is-imported}

{\def\LTcaptype{none} % do not increment counter
\begin{longtable}[]{@{}
  >{\raggedright\arraybackslash}p{(\linewidth - 2\tabcolsep) * \real{0.5000}}
  >{\raggedright\arraybackslash}p{(\linewidth - 2\tabcolsep) * \real{0.5000}}@{}}
\toprule\noalign{}
\begin{minipage}[b]{\linewidth}\raggedright
Component
\end{minipage} & \begin{minipage}[b]{\linewidth}\raggedright
Source
\end{minipage} \\
\midrule\noalign{}
\endhead
\bottomrule\noalign{}
\endlastfoot
Mobility saturation \(M(\rho_\mathrm{max}) = 0\) & Form-FORCED (P4) \\
Specific form \(M(\rho) = M_0(1-\rho/\rho_\mathrm{max})^\beta\) &
Form-FORCED (Universal Mobility Law) \\
Exponent \(\beta\) value & INHERITED per material \\
Stokes-Einstein-class inversion \(\mu \propto 1/M\) & Imported from
kinetic theory \\
Krieger-Dougherty divergence
\(\mu \propto (1 - \rho/\rho_\mathrm{max})^{-\beta}\) & Form-FORCED
(combined) \\
\end{longtable}
}

ED's structural contribution is the \emph{unification} of soft-matter
Universal Mobility Law and suspension-rheology Krieger-Dougherty
divergence under a single architectural principle (P4) plus one
Stokes-Einstein-class inversion step.

\begin{center}\rule{0.5\linewidth}{0.5pt}\end{center}

\subsection{4. D2 --- Shear-Rate
Extension}\label{d2-shear-rate-extension}

\subsubsection{4.1 Strain-rate as P4 saturation
argument}\label{strain-rate-as-p4-saturation-argument}

The architectural template extends naturally from density to strain-rate
magnitude

\[\dot\gamma = \sqrt{2 \, S_{ij} S_{ij}}, \qquad S_{ij} = \frac{1}{2}(\partial_i v_j + \partial_j v_i),\]

where \(S_{ij}\) is the symmetric strain-rate tensor. The natural
rotational invariant \(\dot\gamma\) is a positive scalar with dimensions
of inverse time. Three structurally-distinct sub-cases arise from
different choices of \(M(\dot\gamma)\).

\subsubsection{4.2 Discontinuous shear-thickening (DST) and
shear-thickening
fluids}\label{discontinuous-shear-thickening-dst-and-shear-thickening-fluids}

Apply P4 saturation directly to strain rate:

\[M(\dot\gamma) = M_0 \left(1 - \frac{\dot\gamma}{\dot\gamma_\mathrm{max}}\right)^\beta.\]

Under Stokes-Einstein-class inversion (appropriate for
particle-mobility-driven jamming under shear):

\[\mu(\dot\gamma) = \mu_0 \left(1 - \frac{\dot\gamma}{\dot\gamma_\mathrm{max}}\right)^{-\beta}.\]

Viscosity diverges at the critical shear rate
\(\dot\gamma \to \dot\gamma_\mathrm{max}\), corresponding to
\textbf{shear-rate-driven jamming}. This is the canonical Event Density
prediction for discontinuous shear-thickening (DST) and
shear-thickening-fluid (STF) behavior in dense particulate suspensions.

Empirical correspondence: cornstarch + water at moderate volume fraction
(DST onset at \(\dot\gamma_c \sim 1\)--\(10\,\mathrm{s}^{-1}\)); fumed
silica suspensions (DST at
\(\dot\gamma_c \sim 10^2\)--\(10^4\,\mathrm{s}^{-1}\)); shear-thickening
fluids designed for impact protection (tuned
\(\dot\gamma_c \sim 10^3\)--\(10^5\,\mathrm{s}^{-1}\)). The threshold
\(\dot\gamma_c\) is INHERITED across five orders of magnitude depending
on system.

\subsubsection{4.3 Cross / Carreau
shear-thinning}\label{cross-carreau-shear-thinning}

For fluids whose mobility \emph{increases} with shear (or whose
effective viscosity \emph{decreases} --- polymer-melt physics where
shear aligns chains and reduces effective inter-chain coupling), the
appropriate framework uses monotone non-saturating \(M(\dot\gamma)\).
The Cross-class form admits two derivation paths:

\begin{itemize}
\item
  \textbf{Path A (NS-direct mapping):}
  \(M(\dot\gamma) = M_0/[1 + (\lambda \dot\gamma)^n]\) decreasing;
  viscosity from \(\mu = D \cdot M\):
  \[\mu(\dot\gamma) = \frac{\mu_0}{1 + (\lambda \dot\gamma)^n}.\]
\item
  \textbf{Path B (Stokes-Einstein-class inversion):}
  \(M(\dot\gamma) = M_0[1 + (\lambda \dot\gamma)^n]\) increasing;
  viscosity from \(\mu \propto 1/M\):
  \[\mu(\dot\gamma) = \frac{\mu_0}{1 + (\lambda \dot\gamma)^n}.\]
\end{itemize}

Both paths produce the \textbf{Cross model}. They differ in which
physical mobility ED's \(M\) represents (velocity-mobility for Path A;
particle-mobility for Path B). Path A is the natural physics for
polymer-melt shear-thinning; Path B for suspensions in the
shear-thinning regime.

The Carreau extension adds a high-shear viscosity floor \(\mu_\infty\):

\[\mu(\dot\gamma) = \mu_\infty + (\mu_0 - \mu_\infty)\bigl[1 + (\lambda \dot\gamma)^2\bigr]^{(n-1)/2}.\]

The Carreau-Yasuda generalization adds a shape exponent \(a\). Both are
form-compatible with monotone-\(M\) structures within the architectural
canon; the additional shape parameters \((\mu_\infty, n, a, \lambda)\)
are INHERITED.

Empirical correspondence: polymer melts (HDPE, polystyrene,
polypropylene) with \(n \sim 0.3\)--\(0.5\) and
\(\lambda \sim 10^{-2}\)--\(10^0\,\mathrm{s}\); polymer solutions
(\(n \sim 0.1\)--\(0.4\);
\(\lambda \sim 10^{-1}\)--\(10^1\,\mathrm{s}\)); blood (\(n \sim 0.3\),
\(\lambda \sim 1\)--\(5\,\mathrm{s}\)); paints
(\(n \sim 0.1\)--\(0.4\)).

\subsubsection{\texorpdfstring{4.4 Mixed regimes via non-monotone
\(M(\dot\gamma)\)}{4.4 Mixed regimes via non-monotone M(\textbackslash dot\textbackslash gamma)}}\label{mixed-regimes-via-non-monotone-mdotgamma}

Real concentrated suspensions often exhibit shear-thinning at low shear
rates (alignment / lubrication-mediated) followed by shear-thickening at
higher rates (DST onset). Event Density accommodates this via
non-monotone mobility:

\[M_\mathrm{mixed}(\dot\gamma) = M_0 \bigl[1 + (\lambda_1 \dot\gamma)^{n_1}\bigr] \cdot \left(1 - \frac{\dot\gamma}{\dot\gamma_\mathrm{max}}\right)^\beta.\]

At low \(\dot\gamma\) the shear-thinning factor dominates and \(M\)
rises; near \(\dot\gamma_\mathrm{max}\) the saturation factor drives
\(M \to 0\) and viscosity diverges. The combined curve transitions from
shear-thinning to DST onset.

Empirical correspondence: concentrated colloidal suspensions,
semi-dilute polymer solutions with associative interactions, granular
pastes, food rheology (yogurt, ketchup, mayonnaise).

\subsubsection{4.5 The dual-mapping
framework}\label{the-dual-mapping-framework}

The two distinct mobility-viscosity mappings (Stokes-Einstein-class
inversion for particle-mobility; NS-direct for velocity-mobility) are
not in competition; they apply to different physical scenarios. The
architectural canon supplies the saturation \emph{form}; the choice of
which mobility the form applies to is determined by the system being
modeled. Both choices preserve the form-FORCED / value-INHERITED
methodology cleanly.

\begin{center}\rule{0.5\linewidth}{0.5pt}\end{center}

\subsection{5. D3 --- Rheology-Class
Catalogue}\label{d3-rheology-class-catalogue}

We classify thirteen standard rheology classes by their architectural
status under the P4 + V5 framework. Notation: SE = Stokes-Einstein-class
inversion; NS-direct = velocity-mobility direct mapping; Form-FORCED =
derivable from architectural principles; Compatible-INHERITED =
architecturally consistent but with shape parameters going beyond basic
P4 + V5; Non-canonical = requires structural content beyond P4 + V5.

{\def\LTcaptype{none} % do not increment counter
\begin{longtable}[]{@{}
  >{\raggedright\arraybackslash}p{(\linewidth - 10\tabcolsep) * \real{0.1667}}
  >{\raggedright\arraybackslash}p{(\linewidth - 10\tabcolsep) * \real{0.1667}}
  >{\raggedright\arraybackslash}p{(\linewidth - 10\tabcolsep) * \real{0.1667}}
  >{\raggedright\arraybackslash}p{(\linewidth - 10\tabcolsep) * \real{0.1667}}
  >{\raggedright\arraybackslash}p{(\linewidth - 10\tabcolsep) * \real{0.1667}}
  >{\raggedright\arraybackslash}p{(\linewidth - 10\tabcolsep) * \real{0.1667}}@{}}
\toprule\noalign{}
\begin{minipage}[b]{\linewidth}\raggedright
\#
\end{minipage} & \begin{minipage}[b]{\linewidth}\raggedright
Class
\end{minipage} & \begin{minipage}[b]{\linewidth}\raggedright
Standard form
\end{minipage} & \begin{minipage}[b]{\linewidth}\raggedright
Mapping
\end{minipage} & \begin{minipage}[b]{\linewidth}\raggedright
Status
\end{minipage} & \begin{minipage}[b]{\linewidth}\raggedright
Notes
\end{minipage} \\
\midrule\noalign{}
\endhead
\bottomrule\noalign{}
\endlastfoot
1 & Newtonian & \(\mu = \mathrm{const}\) & Either & Form-FORCED (limit)
& P4 far from saturation \\
2 & Krieger-Dougherty &
\(\mu \propto (1-\phi/\phi_\mathrm{max})^{-[\eta]\phi_\mathrm{max}}\) &
SE & Form-FORCED & Section 3 \\
3 & DST / STF & \(\mu \propto (1-\dot\gamma/\dot\gamma_c)^{-\beta}\) &
SE & Form-FORCED & Section 4.2 \\
4 & Power-law thinning & \(\mu = K\dot\gamma^{n-1}\) & Either &
Compatible-INHERITED & Cross-class high-shear limit \\
5 & Cross & \(\mu = \mu_0/[1+(\lambda\dot\gamma)^n]\) & Either &
Form-FORCED & Section 4.3 \\
6 & Carreau &
\((\mu - \mu_\infty)/(\mu_0-\mu_\infty) = [1+(\lambda\dot\gamma)^2]^{(n-1)/2}\)
& Either & Compatible-INHERITED & Adds \(\mu_\infty\) \\
7 & Carreau-Yasuda & Carreau with shape exponent \(a\) & Either &
Compatible-INHERITED & Industrial standard fit \\
8 & Mixed (thin-then-thick) & Non-monotone \(\mu(\dot\gamma)\) & Either
& Form-FORCED & Section 4.4 \\
9 & Maxwell viscoelastic & \(\tau_R \dot\sigma + \sigma = 2\mu S\) & V5
kernel & Compatible-INHERITED & Section 7 \\
10 & Oldroyd-B & Maxwell + convected derivative + retardation & Beyond
V5 & Non-canonical & Conformation tensor \\
11 & Giesekus & Oldroyd-B + quadratic stress & Beyond V5 & Non-canonical
& Stress nonlinearity \\
12 & FENE-P & Maxwell + finite extensibility & Beyond V5 & Non-canonical
& Polymer chain constraint \\
13 & Bingham yield-stress &
\(\sigma = \sigma_\mathrm{yield} + \mu\dot\gamma\) & Inverted P4 &
Non-canonical & Sign-inverted threshold \\
\end{longtable}
}

\textbf{Summary statistics:}

\begin{itemize}
\tightlist
\item
  \textbf{Form-FORCED (5 classes):} Newtonian, Krieger-Dougherty,
  DST/STF, Cross-class shear-thinning, mixed-regime non-monotone.
\item
  \textbf{Compatible-INHERITED (4 classes):} power-law thinning,
  Carreau, Carreau-Yasuda, Maxwell viscoelastic.
\item
  \textbf{Non-canonical (4 classes):} Oldroyd-B, Giesekus, FENE-P,
  Bingham yield-stress.
\end{itemize}

\textbf{Coverage:} 9 of 13 classes fall within ED's canonical
architectural reach. The four non-canonical classes are concentrated in
polymer-solution-specific models (Oldroyd-B, Giesekus, FENE-P) and
yield-stress fluids (Bingham, requiring sign-inverted P4); each is a
candidate for future architectural extension via tensor /
inverted-saturation primitives, tracked in {[}Candidate Architectural
Extensions 2026{]}.

\begin{center}\rule{0.5\linewidth}{0.5pt}\end{center}

\subsection{6. D4 --- Empirical
Comparison}\label{d4-empirical-comparison}

\subsubsection{6.1 The Universal Mobility Law's ten-system
distribution}\label{the-universal-mobility-laws-ten-system-distribution}

The empirical anchor for the fluid-mechanical extension is the Universal
Mobility Law's ten-system distribution. Table 1 reports the fitted
exponents \(\beta\) and goodness-of-fit \(R^2\) across the ten
chemically distinct materials.

\textbf{Table 1: Universal Mobility Law fits across 10 materials}
{[}Proxmire 2026, Table 1{]}

{\def\LTcaptype{none} % do not increment counter
\begin{longtable}[]{@{}
  >{\raggedright\arraybackslash}p{(\linewidth - 10\tabcolsep) * \real{0.1667}}
  >{\raggedright\arraybackslash}p{(\linewidth - 10\tabcolsep) * \real{0.1667}}
  >{\raggedright\arraybackslash}p{(\linewidth - 10\tabcolsep) * \real{0.1667}}
  >{\raggedright\arraybackslash}p{(\linewidth - 10\tabcolsep) * \real{0.1667}}
  >{\raggedright\arraybackslash}p{(\linewidth - 10\tabcolsep) * \real{0.1667}}
  >{\raggedright\arraybackslash}p{(\linewidth - 10\tabcolsep) * \real{0.1667}}@{}}
\toprule\noalign{}
\begin{minipage}[b]{\linewidth}\raggedright
\#
\end{minipage} & \begin{minipage}[b]{\linewidth}\raggedright
Material
\end{minipage} & \begin{minipage}[b]{\linewidth}\raggedright
Type
\end{minipage} & \begin{minipage}[b]{\linewidth}\raggedright
Mechanism
\end{minipage} & \begin{minipage}[b]{\linewidth}\raggedright
\(\beta\)
\end{minipage} & \begin{minipage}[b]{\linewidth}\raggedright
\(R^2\)
\end{minipage} \\
\midrule\noalign{}
\endhead
\bottomrule\noalign{}
\endlastfoot
1 & Hard-sphere colloids & Colloid & Steric jamming & 1.69 & 0.995 \\
2 & Sucrose-water & Molecular & H-bond viscosity & 2.49 & 0.987 \\
3 & BSA protein & Protein & Hydrodynamic crowding & 2.12 & 0.986 \\
4 & Lysozyme & Protein & Short-range attraction + crowding & 1.36 &
0.998 \\
5 & PMMA colloids & Colloid & Steric jamming & 1.81 & 0.994 \\
6 & Ludox silica & Charged colloid & Electrostatic + steric & 1.41 &
0.999 \\
7 & PEG-water & Polymer & Entanglement & 1.30 & 0.996 \\
8 & Dextran & Polysaccharide & Polymer crowding & 1.46 & 0.993 \\
9 & Casein micelles & Bio-colloid & Depletion + steric & 1.79 & 0.998 \\
10 & Glycerol-water & Small molecule & Viscosity divergence & 1.74 &
0.999 \\
\end{longtable}
}

\textbf{Key statistics:}

\begin{itemize}
\tightlist
\item
  \textbf{Mean} \(\bar\beta = 1.72 \pm 0.37\) (\(N = 10\))
\item
  \textbf{Range} \([1.30, 2.49]\)
\item
  \textbf{Goodness of fit} \(R^2 > 0.986\) in every system; 7 of 10 with
  \(R^2 \geq 0.99\)
\item
  \textbf{Mechanisms covered} 8 distinct physical mechanisms
\item
  \textbf{Materials excluded} 0
\end{itemize}

\subsubsection{\texorpdfstring{6.2 Comparison to canonical
\(\beta = 2.0\)}{6.2 Comparison to canonical \textbackslash beta = 2.0}}\label{comparison-to-canonical-beta-2.0}

The canonical Event Density prediction is \(\beta = 2.0\), arising from
the architectural multiplicative-cooperativity reading of P4 mobility
saturation. The empirical mean \(\bar\beta = 1.72\) deviates from this
central value by:

\[\Delta \beta = |\beta_\mathrm{ED} - \bar\beta| = 0.28.\]

Normalized to the empirical standard deviation:

\[\frac{\Delta \beta}{\sigma_\beta} = \frac{0.28}{0.37} \approx 0.76 \, \sigma.\]

\textbf{The canonical ED prediction is consistent within \(1\sigma\) of
the individual-system scatter level.}

At the standard-error-of-mean level:

\[\sigma_{\bar\beta} = \frac{\sigma_\beta}{\sqrt{N}} = \frac{0.37}{\sqrt{10}} \approx 0.117,\]

\[\frac{\Delta \beta}{\sigma_{\bar\beta}} = \frac{0.28}{0.117} \approx 2.4 \, \sigma.\]

The standard-error-of-mean deviation is \(2.4\sigma\) --- marginal but
admitting structural interpretation: if the ten-system population
under-samples cooperative networks (which would shift \(\bar\beta\)
above 2), the population mean would naturally sit below the canonical
central value. Larger or more diverse system populations would tighten
this comparison.

\subsubsection{\texorpdfstring{6.3 Mechanism-clustered \(\beta\)
patterns}{6.3 Mechanism-clustered \textbackslash beta patterns}}\label{mechanism-clustered-beta-patterns}

Grouping the ten systems by structural cooperativity class:

{\def\LTcaptype{none} % do not increment counter
\begin{longtable}[]{@{}
  >{\raggedright\arraybackslash}p{(\linewidth - 6\tabcolsep) * \real{0.2500}}
  >{\raggedright\arraybackslash}p{(\linewidth - 6\tabcolsep) * \real{0.2500}}
  >{\raggedright\arraybackslash}p{(\linewidth - 6\tabcolsep) * \real{0.2500}}
  >{\raggedright\arraybackslash}p{(\linewidth - 6\tabcolsep) * \real{0.2500}}@{}}
\toprule\noalign{}
\begin{minipage}[b]{\linewidth}\raggedright
Cluster
\end{minipage} & \begin{minipage}[b]{\linewidth}\raggedright
Mechanisms
\end{minipage} & \begin{minipage}[b]{\linewidth}\raggedright
Systems
\end{minipage} & \begin{minipage}[b]{\linewidth}\raggedright
\(\bar\beta\)
\end{minipage} \\
\midrule\noalign{}
\endhead
\bottomrule\noalign{}
\endlastfoot
Cooperative (high \(\beta\)) & H-bond viscosity, hydrodynamic crowding &
sucrose, BSA & \(\bar\beta \approx 2.31\) \\
Steric / canonical (mid \(\beta\)) & Hard-sphere, PMMA, casein, glycerol
& hard-sphere, PMMA, casein, glycerol & \(\bar\beta \approx 1.76\) \\
Gradual / less-cooperative (low \(\beta\)) & Polymer entanglement,
electrostatic, short-range attraction & PEG, lysozyme, Ludox, dextran &
\(\bar\beta \approx 1.38\) \\
\end{longtable}
}

The three-cluster pattern is structurally interpretable under Event
Density's cooperativity reading: higher \(\beta\) corresponds to more
cooperative arrest (more neighbours must be simultaneously displaced for
a particle to move); lower \(\beta\) to more gradual decline. The steric
/ canonical cluster's mean \(\bar\beta \approx 1.76\) sits closest to
the canonical prediction; cooperative networks (H-bond, hydrodynamic)
push above; gradual / less-cooperative systems (polymer entanglement,
electrostatic) pull below.

This pattern is not random material-specific scatter; it reflects a
structurally-coherent gradation of cooperativity strength across the ten
systems.

\subsubsection{6.4 Cross-class extrapolation to broader rheology
literature}\label{cross-class-extrapolation-to-broader-rheology-literature}

Beyond the ten UMD systems, the fluid-mechanical extension predicts that
Krieger-Dougherty divergence, DST exponents, and Cross-class shape
parameters should also fall within form-FORCED bounds, with INHERITED
specific values per material.

\textbf{Krieger-Dougherty \(\beta\) values} in standard
suspension-rheology literature: - Hard-sphere suspensions:
\(\beta \approx 1.5\)--\(2.0\) (canonical
\([\eta]\phi_\mathrm{max} \approx 1.82\)). - Soft / polymeric colloids:
\(\beta \approx 1.0\)--\(2.5\). - Polydisperse suspensions:
\(\beta \approx 1.5\)--\(3.0\) (rises with polydispersity). - Granular
pastes: \(\beta \approx 2\)--\(4\) (frictional contacts). - Fiber / rod
suspensions: \(\beta \approx 2\)--\(5\) (anisotropic alignment).

The wider literature ranges sit at and above the UMD soft-matter range,
consistent with the cooperativity interpretation: granular systems have
stronger frictional cooperativity than colloids; anisotropic systems
have stronger alignment cooperativity than isotropic.

\textbf{DST critical shear rates} span 5+ orders of magnitude across
material classes (cornstarch + water
\(\sim 1\)--\(10\,\mathrm{s}^{-1}\); STF body-armor designs
\(\sim 10^3\)--\(10^5\,\mathrm{s}^{-1}\)). Form-universality holds;
specific \(\dot\gamma_c\) INHERITED.

\textbf{Cross / Carreau shape parameters} \(n \sim 0.1\)--\(0.7\) and
\(\lambda \sim 10^{-3}\)--\(10^1\,\mathrm{s}\) across polymer melts,
solutions, and biological fluids. Form-universality holds; specific
\((n, \lambda)\) INHERITED.

\subsubsection{6.5 Empirical-comparison aggregate
verdict}\label{empirical-comparison-aggregate-verdict}

\begin{itemize}
\tightlist
\item
  \textbf{Form-FORCED universality:} strongly supported. The Universal
  Mobility Law form fits 10 mechanistically distinct systems with
  \(R^2 > 0.986\); cross-class extrapolation to Krieger-Dougherty / DST
  / Cross literature is consistent at order-of-magnitude level.
\item
  \textbf{Canonical \(\beta = 2.0\) quantitative anchor:} consistent
  within \(1\sigma\) at individual-system scatter level. The systematic
  offset (mean 1.72 vs.~canonical 2.0) admits structural interpretation
  via mechanism-class population.
\item
  \textbf{Mechanism-clustered \(\beta\) patterns:} structurally coherent
  with cooperativity interpretation. The three-cluster pattern
  (cooperative / canonical / gradual) is not random scatter.
\end{itemize}

\textbf{Aggregate:} Event Density's fluid-mechanical extension is
consistent with empirical rheology data at qualitative +
light-quantitative level. The form-FORCED structural claim is
empirically validated; the canonical quantitative anchor is consistent
within statistical scatter; mechanism-class patterns add structural
coherence beyond the basic universality finding.

\begin{center}\rule{0.5\linewidth}{0.5pt}\end{center}

\subsection{7. D5 --- Maxwell Viscoelasticity from
V5}\label{d5-maxwell-viscoelasticity-from-v5}

\subsubsection{7.1 V5 cross-chain memory
kernel}\label{v5-cross-chain-memory-kernel-1}

Per arc-N N.2 of the Event Density program, V5 establishes that
same-sector chains acquire correlated bandwidth content via
vacuum-mediated coupling. In fluid-mechanical context, this corresponds
to a finite-memory stress-strain coupling:

\[\sigma_\mathrm{V5}(t) = 2\mu \int_{-\infty}^t K_\mathrm{V5}(t - t') \, S(t') \, dt',\]

where \(K_\mathrm{V5}(t)\) is the V5 memory kernel (zero for \(t < 0\)
by causality; normalized so \(\int_0^\infty K_\mathrm{V5}(t)\,dt = 1\)
by convention), and \(\mu\) is the V5 amplitude prefactor. Both
\(K_\mathrm{V5}\) and \(\mu\) are INHERITED per arc-N N.2 §6.5.

\subsubsection{7.2 Exponential-decay ansatz: the Maxwell-mode
kernel}\label{exponential-decay-ansatz-the-maxwell-mode-kernel}

The simplest single-mode V5 kernel --- and the one most natural under
exponential-decay relaxation kinetics --- is

\[K_\mathrm{V5}(t) = \frac{1}{\tau_R} \, e^{-t/\tau_R} \quad \text{for } t \geq 0,\]

where \(\tau_R\) is the V5 correlation time (INHERITED per material
physics). For polymer melts and solutions
\(\tau_R \in [10^{-3}, 10^0]\,\mathrm{s}\); for biological gels
\(\tau_R \in [10^{-2}, 10^1]\,\mathrm{s}\).

\subsubsection{7.3 Reduction to Maxwell
ODE}\label{reduction-to-maxwell-ode}

Substitute into the V5 stress expression:

\[\sigma(t) = \frac{2\mu}{\tau_R} \int_{-\infty}^t e^{-(t-t')/\tau_R} \, S(t') \, dt'.\]

Differentiate with respect to \(t\):

\[\dot\sigma(t) = \frac{2\mu}{\tau_R} S(t) - \frac{1}{\tau_R} \cdot \frac{2\mu}{\tau_R} \int_{-\infty}^t e^{-(t-t')/\tau_R} \, S(t') \, dt' = \frac{2\mu}{\tau_R} S(t) - \frac{1}{\tau_R} \sigma(t).\]

Rearranging:

\[\boxed{\;\tau_R \dot\sigma + \sigma = 2\mu S.\;}\]

This is the \textbf{Maxwell viscoelastic constitutive equation}
{[}Maxwell 1867; Tanner 2000{]}, the canonical phenomenological model
for linear viscoelasticity. ED's V5 supplies the form-FORCED
memory-kernel structure under exponential-decay ansatz; \(\tau_R\) and
\(\mu\) remain INHERITED.

\subsubsection{7.4 Multi-mode generalized
Maxwell}\label{multi-mode-generalized-maxwell}

For multi-mode viscoelasticity (multiple molecular relaxation processes
--- common in polymer rheology), the V5 kernel becomes a sum of
exponentials:

\[K_\mathrm{V5}^\mathrm{multi}(t) = \sum_i \frac{w_i}{\tau_i} \, e^{-t/\tau_i},\]

with mode weights \(w_i\) and times \(\tau_i\) all INHERITED. This
corresponds to the standard generalized Maxwell model. Form-compatible
with V5; values INHERITED per material structure.

\subsubsection{7.5 V5 vs.~V1 memory: a critical
distinction}\label{v5-vs.-v1-memory-a-critical-distinction}

It is essential to distinguish V5 cross-chain memory from V1
single-chain memory:

\begin{itemize}
\tightlist
\item
  \textbf{V1 memory:} characteristic time
  \(\tau_\mathrm{V1} \sim \ell_P/c \sim 10^{-44}\,\mathrm{s}\) (forced
  by Theorem N1 + Newton-recovery in T19, where
  \(\ell_\mathrm{ED} = \ell_P\)). At any fluid-mechanical scale, V1
  memory coarse-grains to instantaneous response; it is operative only
  at substrate scale.
\item
  \textbf{V5 memory:} characteristic time \(\tau_R\) INHERITED at
  molecular-relaxation scale (\(10^{-3}\)--\(10^0\,\mathrm{s}\) for
  polymer melts, solutions, gels). Operative at fluid-mechanical scales.
\end{itemize}

V5, not V1, is the empirically-relevant viscoelastic memory mechanism at
fluid-mechanical scales. V1's substrate-cutoff regularization is
structurally distinct and is operative only in the smoothness / blow-up
question (Clay-NS territory; out of scope for this paper).

\subsubsection{7.6 Architectural status}\label{architectural-status}

{\def\LTcaptype{none} % do not increment counter
\begin{longtable}[]{@{}
  >{\raggedright\arraybackslash}p{(\linewidth - 2\tabcolsep) * \real{0.5000}}
  >{\raggedright\arraybackslash}p{(\linewidth - 2\tabcolsep) * \real{0.5000}}@{}}
\toprule\noalign{}
\begin{minipage}[b]{\linewidth}\raggedright
Component
\end{minipage} & \begin{minipage}[b]{\linewidth}\raggedright
Source
\end{minipage} \\
\midrule\noalign{}
\endhead
\bottomrule\noalign{}
\endlastfoot
Existence of finite-memory cross-chain stress contribution & Form-FORCED
(V5) \\
Maxwell ODE form \(\tau_R \dot\sigma + \sigma = 2\mu S\) & Form-FORCED
(V5 + Maxwell-mode ansatz) \\
\(\tau_R\) value & INHERITED \\
\(\mu\) amplitude & INHERITED (per arc-N N.2 §6.5) \\
Multi-mode mode-weight distribution & INHERITED \\
Beyond-Maxwell constitutive structure (Oldroyd-B, etc.) & Non-canonical
(requires extensions) \\
\end{longtable}
}

\subsubsection{7.7 Architectural axis structure: P4 +
V5}\label{architectural-axis-structure-p4-v5}

P4 governs \emph{equilibrium / instantaneous} mobility behavior --- what
the steady-state mobility (and viscosity) is at given
\((\rho, \dot\gamma)\). V5 governs \emph{time-dependent memory} --- how
fast stress relaxes after strain-rate changes. Together, P4 and V5 cover
the major architectural axes of non-Newtonian rheology:

{\def\LTcaptype{none} % do not increment counter
\begin{longtable}[]{@{}
  >{\raggedright\arraybackslash}p{(\linewidth - 2\tabcolsep) * \real{0.5000}}
  >{\raggedright\arraybackslash}p{(\linewidth - 2\tabcolsep) * \real{0.5000}}@{}}
\toprule\noalign{}
\begin{minipage}[b]{\linewidth}\raggedright
Behavior axis
\end{minipage} & \begin{minipage}[b]{\linewidth}\raggedright
ED architectural source
\end{minipage} \\
\midrule\noalign{}
\endhead
\bottomrule\noalign{}
\endlastfoot
Equilibrium / instantaneous response & P4 (mobility saturation) \\
Time-dependent memory / relaxation & V5 (cross-chain correlation
kernel) \\
Mixed time-dependent + nonlinear & P4 + V5 combined \\
\end{longtable}
}

\begin{center}\rule{0.5\linewidth}{0.5pt}\end{center}

\subsection{8. Synthesis}\label{synthesis}

\subsubsection{8.1 Three structural
unifications}\label{three-structural-unifications}

The fluid-mechanical extension of the Universal Mobility Law
accomplishes three structurally meaningful unifications:

\begin{enumerate}
\def\labelenumi{\arabic{enumi}.}
\item
  \textbf{Soft-matter mobility ↔ suspension rheology.} The Universal
  Mobility Law (concentration self-diffusion in 10 soft-matter systems
  with \(R^2 > 0.986\)) and the suspension-rheology Krieger-Dougherty
  model (viscosity divergence near close-packing) are revealed as two
  manifestations of a single architectural principle (P4) plus one
  Stokes-Einstein-class inversion step. Two empirical regimes usually
  treated as separate become one Event Density-architectural family.
\item
  \textbf{Equilibrium and time-dependent rheology.} P4 (mobility
  saturation, instantaneous response) and V5 (cross-chain memory kernel,
  Maxwell viscoelasticity) are both architectural canon principles of
  Event Density. Together they cover the major architectural axes of
  non-Newtonian rheology --- instantaneous-mobility classes (Sections
  3--4) and time-dependent-memory classes (Section 7) --- within a
  single framework.
\item
  \textbf{Form-FORCED + canonical-quantitative-anchor methodology.} ED's
  prediction extends beyond pure form-FORCED to include a quantitative
  canonical \(\beta = 2.0\) anchor. The empirical mean
  \(\bar\beta = 1.72 \pm 0.37\) is consistent within \(1\sigma\) of the
  canonical central value at individual-system scatter level. This is
  stronger than form-FORCED-only methodology used in some prior program
  closures and represents a methodological advance: ED can produce
  quantitatively-leaning predictions in addition to structural ones,
  with empirical scatter set by INHERITED material physics.
\end{enumerate}

\subsubsection{8.2 Form-FORCED / value-INHERITED
summary}\label{form-forced-value-inherited-summary}

{\def\LTcaptype{none} % do not increment counter
\begin{longtable}[]{@{}
  >{\raggedright\arraybackslash}p{(\linewidth - 4\tabcolsep) * \real{0.3333}}
  >{\raggedright\arraybackslash}p{(\linewidth - 4\tabcolsep) * \real{0.3333}}
  >{\raggedright\arraybackslash}p{(\linewidth - 4\tabcolsep) * \real{0.3333}}@{}}
\toprule\noalign{}
\begin{minipage}[b]{\linewidth}\raggedright
Component
\end{minipage} & \begin{minipage}[b]{\linewidth}\raggedright
Form status
\end{minipage} & \begin{minipage}[b]{\linewidth}\raggedright
Value status
\end{minipage} \\
\midrule\noalign{}
\endhead
\bottomrule\noalign{}
\endlastfoot
P4 mobility saturation & FORCED & --- \\
Universal Mobility Law functional form & FORCED & --- \\
Krieger-Dougherty divergence
\(\mu \propto (1-\rho/\rho_\mathrm{max})^{-\beta}\) & FORCED & \(\beta\)
INHERITED \\
DST / STF divergence form & FORCED & \(\dot\gamma_\mathrm{max}\),
\(\beta\) INHERITED \\
Cross / Carreau shear-thinning form & FORCED &
\(n, \lambda, \mu_\infty\) INHERITED \\
Mixed-regime non-monotone \(M\) & FORCED via constitutive freedom &
Specific form INHERITED \\
V5 cross-chain memory existence & FORCED-conditional-on-V1 & Amplitude
INHERITED \\
Maxwell ODE form & FORCED via V5 + exponential ansatz & \(\tau_R, \mu\)
INHERITED \\
Canonical \(\beta = 2.0\) central value & Predicted (multiplicative
cooperativity) & --- \\
Empirical \(\bar\beta = 1.72 \pm 0.37\) & --- & INHERITED (population
statistic) \\
\end{longtable}
}

\subsubsection{8.3 Architectural reach}\label{architectural-reach}

Of thirteen catalogued standard rheology classes:

\begin{itemize}
\tightlist
\item
  \textbf{Five form-FORCED:} Newtonian (P4 limit), Krieger-Dougherty,
  DST/STF, Cross-form shear-thinning, mixed-regime non-monotone.
\item
  \textbf{Four compatible-INHERITED:} power-law thinning, Carreau,
  Carreau-Yasuda, Maxwell viscoelastic.
\item
  \textbf{Four non-canonical:} Oldroyd-B, Giesekus, FENE-P, Bingham
  yield-stress.
\end{itemize}

Event Density's architectural canon covers 9 of 13 classes (5 FORCED + 4
INHERITED). The 4 non-canonical classes are concentrated in
polymer-solution-specific models (Oldroyd-B, Giesekus, FENE-P) and
yield-threshold fluids (Bingham); each is a candidate for future
canonical extension via tensor / inverted-saturation primitives.

\begin{center}\rule{0.5\linewidth}{0.5pt}\end{center}

\subsection{9. Discussion}\label{discussion}

\subsubsection{9.1 Significance for Event Density's empirical
reach}\label{significance-for-event-densitys-empirical-reach}

Prior to the fluid-mechanical extension, Event Density's empirical
anchoring rested on the Universal Mobility Law's ten-system validation
in soft matter. The extension developed here moves ED's empirical
content into mainstream non-Newtonian fluid mechanics --- an
empirically-rich domain with extensive industrial, biomedical, and
geophysical applications. The structural unification of soft-matter
mobility, suspension rheology, shear-thickening fluids, polymer-melt
thinning, and Maxwell viscoelasticity under a single architectural
principle (P4 + V5) substantially broadens ED's predictive footprint.

The result is consistent with --- and extends --- the Event Density
program's existing structural-foundations content. The form-FORCED /
value-INHERITED methodology is preserved; the canonical \(\beta = 2.0\)
prediction adds a quantitative anchor that the framework can be tested
against; the four non-canonical classes are honestly catalogued rather
than fitted into the canon by force.

\subsubsection{9.2 Honest limitations}\label{honest-limitations}

Several limitations warrant explicit discussion:

\textbf{Stokes-Einstein-class inversion is imported.} The
mobility-to-viscosity conversion in dense suspensions (Section 3.2) uses
Stokes-Einstein-class physics from standard kinetic theory. Event
Density supplies the saturation form of \(M(\rho)\); kinetic theory
supplies the inversion to viscosity. The combination is
Krieger-Dougherty. ED does not independently derive Stokes-Einstein.
This is honest accounting rather than a weakness --- most empirical
content in physics combines architectural principles with imported
relations.

\textbf{Quantitative comparison is light.} The \(\bar\beta = 1.72\)
vs.~canonical \(\beta = 2.0\) comparison is the strongest quantitative
result. Cross-class extrapolation to broader Krieger-Dougherty / DST /
Cross literature is order-of-magnitude only. Full statistical analysis
(ANOVA on mechanism clusters; formal hypothesis test on canonical
\(\beta\); AIC/BIC against alternative phenomenological models) is
appendix-class work.

\textbf{V5 amplitude INHERITED status limits viscoelastic prediction
power.} Section 7 establishes Maxwell-class compatibility but does not
predict \(\tau_R\) or amplitude values. Real predictive content for
viscoelastic fluids would require additional substrate articulation
beyond the current V5 framework.

\textbf{Standard-error-of-mean vs.~individual-system scatter.} The
canonical \(\beta = 2.0\) prediction is consistent within \(1\sigma\) at
individual-system scatter level (mean \(\bar\beta = 1.72\),
\(\sigma = 0.37\)). At standard-error-of-mean level (which scales as
\(\sigma/\sqrt{N}\)), the deviation is 2.4σ --- marginal but not
strongly inconsistent. Larger or more diverse system populations would
tighten this comparison.

\textbf{Not a Clay-NS resolution.} This paper is ``ED on Earth''
empirical territory --- non-Newtonian rheology --- not Clay-NS
smoothness. The smoothness / blow-up question for the developed-cascade
turbulence regime (closed at Intermediate verdict in NS-3.04) is
structurally distinct from the rheology questions addressed here; the
structural feature breaking BKM-class smoothness in 3D NS is the
advective vortex-stretching term, unrelated to P4 / V5.

\subsubsection{9.3 Non-canonical classes and candidate canon
extensions}\label{non-canonical-classes-and-candidate-canon-extensions}

The four non-canonical classes --- Oldroyd-B, Giesekus, FENE-P, Bingham
yield-stress --- are real empirical phenomena that ED's current
canonical architecture does not natively absorb:

\begin{itemize}
\tightlist
\item
  \textbf{Oldroyd-B} requires a conformation tensor evolving under
  upper-convected-derivative dynamics; this is tensor-extension
  architectural content beyond the scalar P4 + V5 framework.
\item
  \textbf{Giesekus} adds a quadratic stress nonlinearity to Oldroyd-B;
  requires both tensor extension and quadratic-coupling structure.
\item
  \textbf{FENE-P} adds a finite-extensibility constraint on polymer
  chain length; requires polymer-chain-specific structural content.
\item
  \textbf{Bingham yield-stress} requires sign-inverted P4 (mobility zero
  at \emph{lower} stress threshold rather than upper packing limit).
\end{itemize}

Each is a candidate for future architectural extension. Tensor-extension
of the canonical PDE (parallel to the vector-extension framework of
{[}Architectural Canon Vector Extension 2026{]}) could potentially
absorb Oldroyd-B and Giesekus structurally. Inverted-P4 architectural
variants could potentially handle Bingham. None is pursued in this
paper; they are flagged for future canon-extension work.

\subsubsection{9.4 Future work}\label{future-work}

Three concrete future directions emerge from the arc:

\begin{enumerate}
\def\labelenumi{\arabic{enumi}.}
\item
  \textbf{Quantitative empirical refinement.} Full statistical analysis
  of the UMD ten-system distribution (ANOVA on mechanism clusters;
  formal hypothesis test of canonical \(\beta = 2.0\); AIC/BIC
  vs.~alternative models) would strengthen the empirical anchoring
  quantitatively.
\item
  \textbf{Tensor-extension architectural canon.} Following the
  vector-extension framework already articulated for fluid kinematics, a
  tensor-extension would address Oldroyd-B, Giesekus, and FENE-P at the
  architectural level. This is substantial future work parallel in scale
  to a canon-extension memo.
\item
  \textbf{Forced-response and turbulence connections.} A separate arc
  (Reference: {[}P7 Triads NS-Turb Synthesis 2026{]}) found that ED's P7
  triadic harmonic-generation structure provides a restricted-regime
  forced-response prediction in viscous-laminar flow but does not
  architecturally template developed-cascade turbulence. The interplay
  between P4 + V5 (rheology) and P7 (harmonic generation) at the level
  of full nonlinear fluid response remains an open structural question.
\end{enumerate}

\begin{center}\rule{0.5\linewidth}{0.5pt}\end{center}

\subsection{10. Conclusion}\label{conclusion}

We have shown that Event Density's principle P4 (mobility capacity
bound), already empirically validated in the soft-matter regime via the
Universal Mobility Law's ten-system \(R^2 > 0.986\) fit quality, extends
to fluid-mechanical rheology to predict, at the architectural-canon
level:

\begin{itemize}
\tightlist
\item
  \textbf{Krieger-Dougherty viscosity divergence}
  \(\mu(\rho) = \mu_0 (1 - \rho/\rho_\mathrm{max})^{-\beta}\) near
  suspension jamming;
\item
  \textbf{Discontinuous shear-thickening / shear-thickening-fluid (STF)}
  divergence under strain-rate-driven mobility saturation;
\item
  \textbf{Cross / Carreau shear-thinning} under monotone
  strain-rate-dependent mobility;
\item
  \textbf{Mixed thinning-then-thickening regimes} under non-monotone
  constitutive choices;
\item
  \textbf{Maxwell-class viscoelasticity}
  \(\tau_R \dot\sigma + \sigma = 2\mu S\) as the leading-order
  coarse-graining of V5 cross-chain memory kernel under
  exponential-decay ansatz.
\end{itemize}

Of thirteen standard rheology classes catalogued, nine (five
form-FORCED, four compatible-INHERITED) fall within the canonical
architectural reach of P4 + V5; four (Oldroyd-B, Giesekus, FENE-P,
Bingham yield-stress) are non-canonical and flagged for future
architectural extension.

Empirical comparison gives mean \(\bar\beta = 1.72 \pm 0.37\) across the
ten Universal Mobility Law systems, consistent within \(1\sigma\) at
individual-system scatter level with the canonical Event Density
prediction \(\beta = 2.0\). Mechanism-clustered \(\beta\) patterns ---
cooperative networks \(\bar\beta \approx 2.31\), steric / canonical
\(\bar\beta \approx 1.76\), gradual / less-cooperative
\(\bar\beta \approx 1.38\) --- are structurally coherent with the
canonical cooperativity interpretation.

The structural significance is the unification of three
previously-separate empirical regimes --- soft-matter mobility,
suspension-rheology divergence, polymer-class viscoelastic response ---
under a single architectural principle (P4 + V5) with a quantitative
canonical anchor (\(\beta = 2.0\), supported within \(1\sigma\)).
Parameter values (specific \(\beta\) exponent, packing fraction,
shear-rate threshold, Cross / Carreau shape parameters, V5 relaxation
time) remain INHERITED per material-specific physics, consistent with
Event Density's form-FORCED / value-INHERITED methodology.

The result extends Event Density's empirical reach from the soft-matter
regime (Universal Mobility Law) into mainstream non-Newtonian fluid
mechanics, with no new free parameters introduced beyond those already
present in standard rheology models.

\begin{center}\rule{0.5\linewidth}{0.5pt}\end{center}

\subsection{Appendix A. Derivation of the Krieger-Dougherty
Form}\label{appendix-a.-derivation-of-the-krieger-dougherty-form}

We outline the explicit derivation chain for the Krieger-Dougherty
viscosity-divergence form from the architectural principle P4.

\textbf{Step 1 --- P4 mobility saturation.} The architectural principle
P4 specifies that the mobility function \(M(\rho)\) vanishes at a
packing-class upper bound:

\[M(\rho_\mathrm{max}) = 0, \qquad M(\rho) > 0 \text{ for } \rho < \rho_\mathrm{max}.\]

The simplest functional class admissible under non-negativity,
vanishing-at-capacity, smoothness, and dissipative-structure constraints
is the monomial form

\[M(\rho) = M_0 (1 - \rho/\rho_\mathrm{max})^\beta, \qquad \beta > 0.\]

\textbf{Step 2 --- Empirical anchor via Universal Mobility Law.} The
empirical content of P4 in the soft-matter regime is the Universal
Mobility Law \(D(c) = D_0(1 - c/c_\mathrm{max})^\beta\), validated
across 10 chemically distinct systems with \(R^2 > 0.986\) in every fit.

\textbf{Step 3 --- Stokes-Einstein-class inversion.} The dilute-limit
Stokes-Einstein relation

\[D_\mathrm{tracer} = \frac{k_B T}{6\pi \mu a}\]

extends to moderate-to-dense suspensions via the inversion

\[\mu_\mathrm{susp}(\rho) = \mu_0 \cdot \frac{M_0}{M(\rho)},\]

with the proportionality constant set by the dilute limit
(\(\mu_\mathrm{susp}(0) = \mu_0\), the carrier-fluid viscosity).

\textbf{Step 4 --- Combined Krieger-Dougherty form.} Substituting:

\[\mu_\mathrm{susp}(\rho) = \mu_0 (1 - \rho/\rho_\mathrm{max})^{-\beta},\]

structurally identical to the standard Krieger-Dougherty form
\(\mu/\mu_0 = (1 - \phi/\phi_\mathrm{max})^{-[\eta]\phi_\mathrm{max}}\)
under the identification
\(\beta \leftrightarrow [\eta]\phi_\mathrm{max} \approx 1.82\) for hard
spheres.

\textbf{Step 5 --- Form-FORCED / value-INHERITED status.} The functional
form is form-FORCED by P4 + Universal Mobility Law + Stokes-Einstein.
The exponent \(\beta\) value is INHERITED per particle geometry, surface
chemistry, and inter-particle interaction structure.

\begin{center}\rule{0.5\linewidth}{0.5pt}\end{center}

\subsection{Appendix B. The Dual Mobility-Viscosity
Mapping}\label{appendix-b.-the-dual-mobility-viscosity-mapping}

Two distinct mobility-viscosity identifications apply in fluid
mechanics, depending on the physical scenario:

\textbf{Stokes-Einstein-class inversion} (\(\mu \propto 1/M\)). Applies
when \(M\) represents particle self-diffusion in a suspension and
\(\mu\) represents the bulk viscoelastic response. Operative in
dense-suspension Krieger-Dougherty / DST scenarios where particle
mobility drives jamming behavior. Imported from kinetic theory; not
derived from ED.

\textbf{NS-direct mapping} (\(\mu = D \cdot M\)). Applies when \(M\)
represents the kinematic-viscosity-class diffusivity for velocity
gradients in an Event Density vector-extension treatment of fluid
kinematics. Operative in polymer-melt-class shear-thinning scenarios
where chain alignment under shear reduces effective momentum transport.

Both mappings are admissible at the architectural level. The choice is
determined by which physical mobility ED's \(M\) represents in the
particular system --- concentration-mobility (for jamming/DST)
vs.~velocity-mobility (for polymer-melt thinning). The architectural
canon supplies the saturation \emph{form}; the choice of mapping
reflects the physical scenario being modeled. Both choices preserve the
form-FORCED / value-INHERITED methodology cleanly.

\begin{center}\rule{0.5\linewidth}{0.5pt}\end{center}

\subsection{Appendix C. V5 Memory-Kernel Reduction to Maxwell
ODE}\label{appendix-c.-v5-memory-kernel-reduction-to-maxwell-ode}

We derive the Maxwell viscoelastic equation as the leading-order
coarse-graining of V5 cross-chain memory kernel under exponential-decay
ansatz.

\textbf{V5 stress-strain coupling.} Per arc-N N.2, V5 introduces:

\[\sigma_\mathrm{V5}(t) = 2\mu \int_{-\infty}^t K_\mathrm{V5}(t - t') S(t') \, dt'.\]

\textbf{Exponential ansatz.} The simplest single-mode V5 kernel:

\[K_\mathrm{V5}(t) = \frac{1}{\tau_R} e^{-t/\tau_R} \quad (t \geq 0).\]

\textbf{Differentiation.} Taking \(d/dt\) of the V5 stress:

\[\dot\sigma(t) = 2\mu K_\mathrm{V5}(0) S(t) + 2\mu \int_{-\infty}^t \frac{dK_\mathrm{V5}}{dt}(t-t') S(t') \, dt'.\]

For the exponential kernel: - \(K_\mathrm{V5}(0) = 1/\tau_R\). -
\(dK_\mathrm{V5}/dt = -K_\mathrm{V5}/\tau_R\).

Substituting:

\[\dot\sigma(t) = \frac{2\mu}{\tau_R} S(t) - \frac{1}{\tau_R} \cdot 2\mu \int_{-\infty}^t K_\mathrm{V5}(t-t') S(t') \, dt' = \frac{2\mu}{\tau_R} S(t) - \frac{1}{\tau_R} \sigma(t).\]

\textbf{Maxwell ODE.} Rearranging:

\[\tau_R \dot\sigma + \sigma = 2\mu S.\]

\textbf{Multi-mode generalization.} For multiple molecular relaxation
processes:

\[K_\mathrm{V5}^\mathrm{multi}(t) = \sum_i \frac{w_i}{\tau_i} e^{-t/\tau_i},\]

corresponding to the standard generalized Maxwell model with mode
weights \(w_i\) and times \(\tau_i\) all INHERITED.

\textbf{Architectural status.} Form-FORCED at the
existence-of-memory-kernel + Maxwell-mode-ansatz level via V5 +
exponential-decay ansatz; values INHERITED per molecular-relaxation
physics.

\begin{center}\rule{0.5\linewidth}{0.5pt}\end{center}

\subsection{References}\label{references}

{[}Architectural Canon 2026{]} Proxmire, A. \emph{The Architectural
Canon of Event Density.} Event Density program, March 2026.

{[}Born-Gleason 2026{]} Proxmire, A. \emph{Born Rule from Gleason-Busch
in Event Density.} ED Foundational Theorems series, 2026.

{[}Candidate Architectural Extensions 2026{]} \emph{Candidate
Architectural Extensions} tracking memo, Event Density program, 2026.

{[}Foundations 2025{]} Proxmire, A. \emph{Foundations of Event Density:
Axiomatic Derivation and Universality.} Event Density program, 2025.

{[}Krieger \& Dougherty 1959{]} Krieger, I.M. \& Dougherty, T.J. ``A
Mechanism for Non-Newtonian Flow in Suspensions of Rigid Spheres.''
\emph{Trans. Soc. Rheol.} 3, 137--152 (1959).

{[}Maxwell 1867{]} Maxwell, J.C. ``On the dynamical theory of gases.''
\emph{Philos. Trans. R. Soc.} 157, 49--88 (1867).

{[}NS-1.05 2026{]} Proxmire, A. \emph{NS-1.05 --- Synthesis: Final B2
Dimensional-Forcing Verdict.} Event Density Navier-Stokes roadmap, 2026.

{[}NS-3.04 2026{]} Proxmire, A. \emph{NS-3.04 --- Synthesis + Path C
Verdict.} Event Density Navier-Stokes roadmap, 2026.

{[}P7 Triads NS-Turb Synthesis 2026{]} Proxmire, A. \emph{P7 Triads
NS-Turb-5 Synthesis.} Event Density Navier-Stokes Turbulence arc, 2026.

{[}Proxmire 2026{]} Proxmire, A. \emph{Universal Degenerate-Mobility
Scaling in Crowded Soft Matter.} Event Density program / Universal
Mobility Law paper, 2026.

{[}Russel, Saville \& Schowalter 1989{]} Russel, W.B., Saville, D.A. \&
Schowalter, W.R. \emph{Colloidal Dispersions.} Cambridge University
Press, 1989.

{[}Tanner 2000{]} Tanner, R.I. \emph{Engineering Rheology.} 2nd ed.,
Oxford University Press, 2000.

\begin{center}\rule{0.5\linewidth}{0.5pt}\end{center}

\emph{Manuscript draft. Event Density program --- P4 Non-Newtonian arc.
Author: Allen Proxmire. AI collaborator: Claude. April 2026.}

\end{document}
